% !TEX program = lualatex
% Generated by new_lecture_note.sh
\documentclass[11pt]{article}

\usepackage{fontspec}
\usepackage{microtype}
\usepackage{mathtools,amssymb,amsfonts,amsthm,bm}
\usepackage{enumitem}
\usepackage{booktabs,array,xcolor}
\usepackage{geometry,fancyhdr,setspace}
\usepackage[hidelinks]{hyperref}
\usepackage[nameinlink,noabbrev]{cleveref}

\geometry{margin=1in,headheight=15pt}
\setstretch{1.12}
\setlength{\parindent}{0pt}
\setlength{\parskip}{0.55em plus 0.12em minus 0.08em}
\setlength{\abovedisplayskip}{0.8em plus 0.2em minus 0.1em}
\setlength{\belowdisplayskip}{0.8em plus 0.2em minus 0.1em}
\allowdisplaybreaks[3]
\setlist{leftmargin=*,topsep=0.35em,itemsep=0.15em,parsep=0pt}

% ---------- Metadata ----------
\newcommand{\studentname}{Keshav Ramamurthy}
\newcommand{\coursename}{Math 113}
\newcommand{\courselongname}{Abstract Algebra}
\newcommand{\lecturenumber}{04}
\newcommand{\lecturetitle}{Lecture 04}
\newcommand{\lecturedate}{\today}

% ---------- Reusable notation ----------
\newcommand{\N}{\mathbb{N}}
\newcommand{\Z}{\mathbb{Z}}
\newcommand{\Q}{\mathbb{Q}}
\newcommand{\R}{\mathbb{R}}
\newcommand{\C}{\mathbb{C}}
\newcommand{\F}{\mathbb{F}}
\newcommand{\K}{\mathbb{K}}
\newcommand{\E}{\mathbb{E}}
\newcommand{\Prob}{\mathbb{P}}
\newcommand{\1}{\mathbf{1}}
\newcommand{\eps}{\varepsilon}
\newcommand{\dd}{\,\mathrm{d}}
\newcommand{\abs}[1]{\left\lvert#1\right\rvert}
\newcommand{\norm}[1]{\left\lVert#1\right\rVert}
\newcommand{\ip}[2]{\left\langle#1,#2\right\rangle}
\newcommand{\set}[1]{\left\{#1\right\}}
\newcommand{\given}{\,\middle\vert\,}
\newcommand{\indep}{\perp\!\!\!\perp}
\newcommand{\lto}{\longrightarrow}
\newcommand{\deriv}[2]{\frac{\mathrm{d}#1}{\mathrm{d}#2}}
\newcommand{\pderiv}[2]{\frac{\partial#1}{\partial#2}}
\DeclareMathOperator{\Aut}{Aut}
\DeclareMathOperator{\Hom}{Hom}
\DeclareMathOperator{\Ker}{ker}
\DeclareMathOperator{\im}{im}
\DeclareMathOperator{\rank}{rank}
\DeclareMathOperator{\nullity}{nullity}
\DeclareMathOperator{\Span}{span}
\DeclareMathOperator{\Var}{Var}
\DeclareMathOperator{\Cov}{Cov}

% ---------- Note environments ----------
\newtheoremstyle{notestyle}{0.7em}{0.7em}{\itshape}{}{}{\bfseries}{.5em}{}
\theoremstyle{notestyle}
\newtheorem{theorem}{Theorem}[section]
\newtheorem{lemma}[theorem]{Lemma}
\newtheorem{proposition}[theorem]{Proposition}
\newtheorem{corollary}[theorem]{Corollary}
\theoremstyle{definition}
\newtheorem{definition}[theorem]{Definition}
\newtheorem{example}[theorem]{Example}
\newtheorem{remark}[theorem]{Remark}
\newenvironment{claim}{\par\noindent\textbf{Claim.}\ }{\par}
\newenvironment{idea}{\par\noindent\textit{Idea.}\ }{\par}
\newenvironment{proofsketch}{\par\noindent\textbf{Proof sketch.}\ }{\par}

% ---------- Header ----------
\pagestyle{fancy}
\fancyhf{}
\fancyhead[L]{\coursename}
\fancyhead[C]{\lecturetitle}
\fancyhead[R]{\studentname}
\fancyfoot[C]{\thepage}

\title{\vspace{-1.5em}\textbf{\coursename\ -- \lecturetitle}}
\author{\studentname}
\date{\lecturedate}

\begin{document}
\maketitle

\section{Cyclic Groups}
\begin{definition}[Cyclic]
\label{def:label}
A group $G$ is cyclic if $$\exists a \in G \text{ such that } \langle a \rangle
=G(\text{ where } \langle a \rangle \rightarrow \{\cdots, a^{-2}, a^{-1}, a^{0} = e, a^{1}, \cdots\})$$
\end{definition}
\begin{theorem}[Commutativity]
\label{thm:label}
Every Cyclic Group is Abelian.
\end{theorem}
\begin{proof}
    Let $G$ be cyclic and $a \in G$ a generator of g.
\par    Let $b, c \in G$, notice that $a^{x} = b, a^{y} = c$ for $x, y \in \mathbb{Z}$ WLOG assume $x \le y$
Notice that $$b * c = a^n * a^m = a^{n + m}$$ and $$c * b = a^{m} * a^{n} = a^{m + n}
=a^{m + n}$$
So commutativity holds and $G$ is Abelian.
\end{proof}
\begin{theorem}[Division in $\mathbb{Z}$]
\label{thm:label}
If $n \in \mathbb{Z}$ and $m \in \mathbb{Z}^{+}$ then $$\exists! q, r, \in \mathbb{Z}
\text{ such that } n = mq + r \textbf{ and } 0 \le r < m $$
\begin{remark}[$\exists !$]
\label{rem:label}
$\exists !$ denotes "there exists unique"
\end{remark}
\end{theorem}
\begin{theorem}[Subgroup Cyclic Properties]
\label{thm:label}
\textbf{Any subgroup of a cyclic group is cyclic}
\end{theorem}
\begin{proof}
$G$ is cyclic and $\langle a \rangle = G$ and $H \le G.$
\par If $H$ is the identity, we have that $H = \langle e\rangle$ is cyclic.
\par If $H$ is not the identity, it contains $k \ne e$. Since $\langle a \rangle = G$,
we can write $k = a^{n}$ for some $n \in \mathbb{Z}$ with $n \ne 0$ (as $k \ne e$).
\par \textbf{Defining the minimal exponent.} Consider the set of positive exponents
whose powers land in $H$:
$$S = \{j \in \mathbb{Z}^{+} \mid a^{j} \in H\}.$$
$S$ is nonempty: if $n > 0$ then $n \in S$ directly; if $n < 0$ then
$a^{-n} = (a^{n})^{-1} = k^{-1} \in H$, since $H$ is a subgroup (closed under
inverses), so $-n \in S$. By the well-ordering of $\mathbb{Z}^{+}$, $S$ has a least
element
$$m = \min S, \qquad \text{i.e. } a^{m} \in H \text{ but } a^{j} \notin H \text{ for every } 0 < j < m.$$
This is the \textbf{minimality of $m$} — the one fact about $m$ we exploit below.
\par Goal: show $H = \langle a^{m}\rangle$; then $H$ is cyclic. The containment
$\langle a^{m}\rangle \subseteq H$ is free: $a^{m} \in H$ and $H$ is closed under
products and inverses, so every power $(a^{m})^{q}$, $q \in \mathbb{Z}$, stays in $H$
(cyclic-subgroup theorem, Lecture 03).
\par For the other direction, take any $b \in H$. Since $H \subseteq G = \langle a
\rangle$, we have $b = a^{n}$ for some $n \in \mathbb{Z}$, and it suffices to show
$m \mid n$.
\par Use the division property: $n = qm + r$ with $0 \le r < m$, and we want to show
that $r = 0$.
We get $$a^n = (a^m)^q * a^r$$
Then inverse it up:
$$a^{r} = a^{\,n - qm} = (a^{m})^{-q} * a^{n}$$
Both factors lie in $H$ ($a^{n} = b \in H$, and $(a^{m})^{-q} \in H$ by closure under
inverses and products), so by closure $a^{r} \in H$.
\par \textbf{Why $r = 0$ (the minimality step).} We now know $a^{r} \in H$ together
with $0 \le r < m$. Suppose toward a contradiction that $r > 0$. Then $r \in
\mathbb{Z}^{+}$, $r < m$, and $a^{r} \in H$ — that is, $r \in S$ with
$r < m = \min S$, contradicting the minimality of $m$. So $r$ cannot be positive, and
$0 \le r$ then forces $r = 0$.
\par Hence $m \mid n$, so $b = a^{n} = (a^{m})^{q} \in \langle a^{m} \rangle$. Since
$b \in H$ was arbitrary, $H \subseteq \langle a^{m} \rangle$, and therefore
$H = \langle a^{m} \rangle$ is cyclic.
\end{proof}
\begin{corollary}[Scalars]
\label{cor:label}
The subgroups of $\langle \mathbb{Z}, + \rangle$ are the groups $\langle n\mathbb{Z}, +\rangle$
for $n \in \mathbb{Z}_{\ge 0}$
\end{corollary}
\begin{proof}
$\langle \mathbb{Z}, + \rangle$ is cyclic, generated by $1$: in additive notation the
powers $a^{j}$ become multiples $j \cdot 1$, and every $z \in \mathbb{Z}$ is
$z = z \cdot 1$. So the theorem above applies with $a = 1$, reading "$a^{j}$" as
"$j$".
\par Let $H \le \langle \mathbb{Z}, + \rangle$. If $H = \{0\}$, then $H = 0\mathbb{Z}$
and we are done. Otherwise $H$ contains some $k \ne 0$; if $k < 0$ then $-k \in H$
(inverses), so $H$ contains a positive integer and, by well-ordering, a least
positive one:
$$m = \min\{j \in \mathbb{Z}^{+} \mid j \in H\}.$$
\par We claim $H = m\mathbb{Z}$. Since $m \in H$ and $H$ is closed under addition and
negation, $qm \in H$ for every $q \in \mathbb{Z}$, so $m\mathbb{Z} \subseteq H$.
Conversely, let $h \in H$ and divide: $h = qm + r$ with $0 \le r < m$. Then
$$r = h - qm \in H$$
by closure under differences, so the minimality of $m$ forces $r = 0$ exactly as in
the theorem above. Hence $m \mid h$, i.e. $h \in m\mathbb{Z}$, giving
$H = m\mathbb{Z} = \langle m\mathbb{Z}, + \rangle$.
\par Conversely, every $n\mathbb{Z}$ with $n \in \mathbb{Z}_{\ge 0}$ really is a
subgroup: $0 = n \cdot 0 \in n\mathbb{Z}$, and $nx - ny = n(x - y) \in n\mathbb{Z}$.
\par So the subgroups of $\langle \mathbb{Z}, + \rangle$ are exactly the
$\langle n\mathbb{Z}, + \rangle$ for $n \in \mathbb{Z}_{\ge 0}$.
\end{proof}
\begin{theorem}[GCD]
\label{thm:label}
If $r, s \in \mathbb{Z}^{+}$, then $H = \{nr + ms \mid n, m \in \mathbb{Z}\}$ is a
subgroup of $\langle \mathbb{Z}, +\rangle$.
\end{theorem}
\begin{proof}
First, $H \ne \emptyset$: $0 = 0 \cdot r + 0 \cdot s \in H$ (and note
$r = 1 \cdot r + 0 \cdot s \in H$ and $s = 0 \cdot r + 1 \cdot s \in H$).
\par Now check closure under differences (the one-step subgroup test). Take
$h_1 = n_1 r + m_1 s$ and $h_2 = n_2 r + m_2 s$ in $H$, with $n_i, m_i \in
\mathbb{Z}$. Then
$$h_1 - h_2 = (n_1 - n_2) r + (m_1 - m_2) s$$
which is again an integer multiple of $r$ plus an integer multiple of $s$, so
$h_1 - h_2 \in H$. A nonempty subset of a group that is closed under
$x y^{-1}$ (here, under $x - y$) is a subgroup, so $H \le \langle \mathbb{Z},
+\rangle$.
\end{proof}
\begin{remark}[What is H?]
\label{rem:label}
We will in fact have that $H = d\mathbb{Z}$ where $d = \gcd(r, s)$:
\par By the Scalars corollary, $H = n\mathbb{Z}$ for some $n \in \mathbb{Z}_{\ge 0}$.
Since $r \in H$ and $r > 0$, we must have $n \ne 0$; write $d = n \in
\mathbb{Z}^{+}$. Because $r, s \in H = d\mathbb{Z}$, both are multiples of $d$.
\par Also $d \in d\mathbb{Z} = H$, so $d$ itself has the form $d = nr + ms$ for some
$n, m \in \mathbb{Z}$ — a Bézout identity for the gcd.
\end{remark}
\begin{proposition}[What is D?]
\label{prop:label}
$d$ is a divisor of both $r$ and $s$, and any other such divisor $d'$ must divide
$d$ (so $d$ is the greatest common divisor, $d = \gcd(r, s)$).
\end{proposition}
\begin{proof}
We already showed $d \mid r$ and $d \mid s$: both $r$ and $s$ lie in $d\mathbb{Z} = H$.
\par Let $d'$ be any common divisor of $r$ and $s$, so $r = d'x$ and $s = d'y$ for
some $x, y \in \mathbb{Z}$. Feed these into the Bézout form of $d$ from the remark:
$$d = nr + ms = n(d'x) + m(d'y) = d'(nx + my),$$
so $d' \mid d$. Every common divisor of $r$ and $s$ divides $d$; in particular no
common divisor is larger than $d$, which is exactly what makes $d$ the greatest
common divisor.
\end{proof}
\begin{definition}[Order of an element]
\label{def:element-order}
The order of $a \in G$, written $\operatorname{ord}(a)$, is the least
$n \in \mathbb{Z}^{+}$ with $a^{n} = e$, if such an $n$ exists; otherwise $a$ has
infinite order.
\par The "least" is the same minimality move as before: if $a^{k} = e$ for any
$k \ne 0$, then also $a^{-k} = e$, so $\{j \in \mathbb{Z}^{+} \mid a^{j} = e\}$ is
nonempty and well-ordering produces the least element.
\end{definition}
\begin{theorem}[Classification of Cyclic Groups]
\label{thm:classification}
Let $G = \langle a \rangle$ be cyclic. If $a$ has infinite order, then
$G \cong \langle \mathbb{Z}, + \rangle$; if $\operatorname{ord}(a) = n < \infty$,
then $G \cong \langle \mathbb{Z}_n, + \rangle$.
\end{theorem}
\begin{proof}
\textbf{Case 1: infinite order.} Define
$$\varphi : \mathbb{Z} \to G, \qquad \varphi(k) = a^{k}.$$
\par Homomorphism: $\varphi(j + k) = a^{\,j + k} = a^{j} * a^{k} =
\varphi(j) * \varphi(k)$ (exponent law).
\par Surjective: every element of $G = \langle a \rangle$ is some $a^{k} = \varphi(k)$.
\par Injective: suppose $\varphi(j) = \varphi(k)$ with $j \ne k$, WLOG $j > k$. Then
$$a^{\,j - k} = a^{j} * (a^{k})^{-1} = e$$
with $j - k \in \mathbb{Z}^{+}$, contradicting infinite order (no positive power of
$a$ is $e$). So $\varphi$ is an isomorphism and $G \cong \langle \mathbb{Z},
+ \rangle$.
\par \textbf{Case 2: $\operatorname{ord}(a) = n < \infty$.} The key fact is
$$a^{i} = a^{j} \iff i \equiv j \ (\mathrm{mod}\ n).$$
\par ($\Leftarrow$): if $i = j + qn$, then $a^{i} = a^{j} * (a^{n})^{q} = a^{j} * e
= a^{j}$.
\par ($\Rightarrow$): if $a^{i} = a^{j}$, then $a^{\,i - j} = a^{i} * (a^{j})^{-1} =
e$. Divide (the division property applies to any integer):
$i - j = qn + r$ with $0 \le r < n$. Then
$$a^{r} = a^{\,i - j} * (a^{n})^{-q} = e,$$
and minimality of $n$ (the least positive exponent with power $e$) forces $r = 0$,
exactly as in the subgroup proof. So $n \mid i - j$, i.e. $i \equiv j \pmod n$.
\par Now define
$$\varphi : \mathbb{Z}_n \to G, \qquad \varphi([k]) = a^{k}.$$
The thing to check for a map out of $\mathbb{Z}_n$ is that it is well-defined:
if $[k] = [k']$ then $k \equiv k' \pmod n$, so $a^{k} = a^{k'}$ by
($\Leftarrow$), and $\varphi$ does not depend on the representative.
\par Homomorphism: $\varphi([j] + [k]) = \varphi([j + k]) = a^{\,j + k} = a^{j} *
a^{k} = \varphi([j]) * \varphi([k])$.
\par Surjective as before, and injective by ($\Rightarrow$):
$\varphi([j]) = \varphi([k])$ gives $j \equiv k \pmod n$, i.e. $[j] = [k]$.
\par So $G \cong \langle \mathbb{Z}_n, + \rangle$; in particular the distinct
elements of $G$ are exactly $e, a, \dots, a^{n - 1}$, so $|G| = n =
\operatorname{ord}(a)$.
\end{proof}






\section{Definitions and notation}
\begin{definition}[First definition]
Write the definition and one sentence explaining why it matters.
\end{definition}

\section{Claims, theorems, and proof sketches}
\begin{claim}
State the claim in one precise sentence.
\end{claim}

\begin{proofsketch}
Record the key idea, the first useful equation, and the step that still needs checking.
\end{proofsketch}

\section{Examples and calculations}
\begin{example}[Lecture example]
Write the setup, the calculation, and the conclusion.
\end{example}

\section{Questions and follow-up}
\begin{itemize}
  \item What is still unclear?
  \item Which exercise or reading should I revisit?
  \item TODO:
\end{itemize}

\end{document}
